\scanpage{88}
$(\mathrm{iv}^{(\mathrm{iv})})$ $\sigma_{0V}$ est une réflexion, ou l'identité
(avec \deleted{\unsure{la|sa}} remarque)

\uline{NB} Dans le cas impropre, $\sigma_{0\unsure{E|k}}$ est la sym. p.r. à
un pt (car. $\ne2$), $\sigma_{0V}$ est la symétrie $\deleted{\unsure{v}}v\mapsto-v$.

\uline{Points fixes sous} $G$ = \uline{pts fixes sous} $G^+$ = \uline{pts}
\uline{fixes sous} $u$ :
% slanted
\margin{\uline{terminologie} \uline{Centre} de la réalisation géom de $\Pi_\infty$\ldots}
\[ x_0=y_0=\frac{1}{2-\alpha} \tag*{(29)} \]
(qui existe ssi $\alpha\ne2$, et est tjrs unique)

4. \uline{Valeurs propres de} $u_V$. \uline{Comparaison avec}
\uline{les} $\unsure{\varphi\zeta|\varphi_\zeta}$\ldots

\[ u_V=\begin{pmatrix}\alpha&-1\\1&0\end{pmatrix} \tag*{(30)} \]
\[ \deleted{\unsure{\det(|\det}} \left\{\begin{aligned}\det u_V&=1\\
\operatorname{Tr} u_V&=\alpha\end{aligned}\right. \tag*{(31)} \]
\unsure{équations|équation} aux valeurs propres $(\zeta,\in\overline{k})$
\[ \unsure{\zeta^2|\zeta}-\alpha\zeta+1=0\qquad\text{discriminant} \tag*{(32)} \]
\[ \Delta\unsure{(\alpha)|{}}=\alpha^2-4=(\alpha-2)(\alpha+2) \tag*{(33)} \]
En termes du choix d'\uline{un} tel $\zeta$, on
trouve
\[ \alpha=\zeta+\zeta^{-1} \tag*{(33)} \]
( \uline{NB} les deux racines $\zeta,\zeta'$ de \unsure{notre|cette} équation
satisfont
\[ \left\{\begin{aligned}\zeta+\zeta'&=\alpha\\
\zeta\zeta'&=1\end{aligned}\right.\qquad\text{i.e. }\zeta'=\zeta^{-1}. \tag*{(34)} \]
% archive sheet number: 87
